\section{Related work}
\label{sec:related}

We organise the related work by the question that each line of work answers. Each subsection ends with the
difference from the present work. Table~\ref{tab:positioning} collects the closest works, and the paragraph after
it states the novelty claim.

\subsection{Lifelong behaviour models}
\label{sec:rel-lifelong}

Lifelong behaviour models in click-through prediction retrieve information from long user histories. SIM
\citep{pi2020sim} selects behaviours from a lifelong sequence that match the category of the target item, and passes
them to a deep CTR model. TWIN \citep{chang2023twin} applies a two-stage target attention over a lifelong sequence on
the Kuaishou platform, which belongs to the same platform family as KuaiRand. Both learn from policy logs and rely on
item identifiers that the history shares with the target.
\textit{Difference.} LEAF uses exposure-conditional statistics in which non-clicks enter the denominator, transfers to
a pool that is disjoint from the lifelong history through content keys only, and is evaluated on randomised insertions
under a temporal protocol.

\subsection{Click models and examination}
\label{sec:rel-click}

A click model explains a click as the product of an examination event and a judgement about the item. The examination
hypothesis and its position-based form \citep{richardson2007predicting} make the examination probability depend on the
position of the item. \citet{craswell2008cascade} compare cascade and position-bias models on search logs, and the
monograph of \citet{chuklin2015clickmodels} summarises this family. The dynamic Bayesian network model
\citep{chapelle2009dbn} separates perceived relevance from actual relevance through a satisfaction state.
\textit{Difference.} DAX conditions examination on the user, through a lifelong prior, and on the scenario, through a tab
factor; it has no position term and no satisfaction state, and it targets item appeal under uniform insertion, which
we check on randomised exposure.

\subsection{Exposure-aware recommendation}
\label{sec:rel-exposure}

A line of work learns from impressions and not only from clicks. Impression-aware recommenders treat shown but unclicked
items as evidence and use the impression set as features or as negative samples \citep{perezmaurera2025iars}. ExpoMF
\citep{liang2016exposure} models exposure as a latent variable that is fitted by EM. Inverse-propensity estimators for
missing-not-at-random implicit feedback \citep{saito2020unbiased} weight the loss by estimated exposure propensities, and
\citet{joachims2017unbiased} apply the same idea to clicks with examination propensities. \citet{wang2018regressionem}
and \citet{wang2016selection} estimate position bias and relevance jointly, without randomisation, in personal search.
ESMM \citep{ma2018esmm} models the entire impression space, and \citet{krause2024coexposure} use multinomial-logit
choice with item co-exposure. cEASE \citep{jeunen2020cease} adds item side information to a closed-form model.
\textit{Difference.} The impression channel alone is not claimed as new here, and DAX differs in that it is a generative
model of examination and appeal fitted by EM, whereas the estimators above weight a loss or use impressions as negatives
or features.

\subsection{Occupancy, detection and small-area estimation}
\label{sec:rel-occupancy}

Occupancy models separate the presence of a species from its detection in repeated surveys \citep{mackenzie2002occupancy}:
a site can be occupied and still go undetected, and the pattern of repeated non-detections identifies both quantities.
Small-area estimation shrinks noisy direct estimates for many small areas towards a regression on area covariates
\citep{fay1979smallarea}, and multilevel regression with poststratification estimates population quantities from
non-representative samples \citep{park2004mrp}. We found no recommendation paper that uses the occupancy likelihood to
estimate item appeal from a policy log.
\textit{Difference.} The likelihood of DAX has the same structure as the occupancy model, with an impression playing the
role of a survey and a click the role of a detection, but its estimand is item appeal under a uniform audience, and the
borrowed MRP and Fay-Herriot forms, which we tested on this target, are weaker than the multiplicative form.

\subsection{Debiasing for graph recommenders}
\label{sec:rel-graph}

Graph collaborative filtering propagates signals over the user-item graph, and bias corrections have been applied to that
propagation. NAVIP \citep{kim2022navip} debiases neighbour aggregation in a graph neural network for recommendation.
Popularity-bias work \citep{zhang2021pda} applies a causal intervention to the popularity effect. AutoDebias
\citep{chen2021autodebias} learns debiasing parameters from a small uniform sample, which is not admissible under our
protocol because the randomly exposed data serve only for validation and test. Two further graph methods, adversarial
edge dropout and the removal of a popularity direction from the embeddings after training, are relevant here
\citep{zhang2024advdrop,islam2026ppd}; neither is re-run in this paper.
The backbones used in our experiments are LightGCN \citep{he2020lightgcn}, EASE \citep{steck2019ease} and the
r-AdjNorm normalisation \citep{zhao2022radjnorm}.
\textit{Difference.} The graph methods above act on propagation or on the embeddings and use no impressions and no
lifelong data, whereas DAX re-estimates item appeal from the impressions under a different audience, and LEAF adds
user-side behaviour on other items.

\subsection{Evaluation under randomised exposure}
\label{sec:rel-eval}

Random insertion is a direct test of how a recommender responds to items that the policy did not choose. KuaiRand
\citep{gao2022kuairand} ships randomly exposed videos in its feed logs, and KuaiRec \citep{gao2022kuairec} is a fully
observed matrix that supports the evaluation of recommenders. Several recent papers use the random insertions of KuaiRand
for evaluation: propensity calibration \citep{hu2025uncertainty}, debiasing of watch-time signals \citep{zhao2023d2co},
an analysis of whether the evaluation of debiasing methods is itself biased \citep{wang2025debiasunreliable}, and an audit
of category means in logs with random insertions \citep{zhang2026categorymeans}. Classical debiasing benchmarks use a
small randomly exposed sample from Coat and Yahoo! R3 \citep{schnabel2016, marlin2009collaborative}.
\textit{Difference.} As far as our searches went, none of these works uses the out-of-pool history of each user in the
random test, which is the source of the LEAF channels here.

\subsection{Positioning and novelty}
\label{sec:rel-position}

\begin{table}[t]
\caption{Closest works and the difference from the present work.}
\label{tab:positioning}
\centering
\footnotesize
\raggedright
\begin{tabular}{p{3.4cm}p{4.2cm}p{5.4cm}}
\toprule
Work & What it estimates, and from which data & Difference from this work \\
\midrule
SIM \citep{pi2020sim}; TWIN \citep{chang2023twin} & CTR from lifelong sequences; policy logs; history and target share item ids & Exposure-conditional statistics; transfer to a disjoint pool through content keys; randomised test; SIM and TWIN not run on the main setting \\
Examination models \citep{richardson2007predicting,craswell2008cascade,chuklin2015clickmodels} & Click as examination times attractiveness; position or list scan & User and scenario examination; no position term; uniform-insertion target \\
Dynamic Bayesian network \citep{chapelle2009dbn} & Perceived versus actual relevance with satisfaction & No satisfaction state \\
Regression-EM \citep{wang2018regressionem,wang2016selection} & Position bias and relevance without randomisation; personal search & Item appeal under a uniform audience; generative fit by EM over user and tab \\
ExpoMF \citep{liang2016exposure} & Latent exposure fitted by EM & Impressions are observed; latent examination given an impression \\
Unbiased LTR \citep{joachims2017unbiased} & Inverse-propensity weighted losses & Generative model instead of a weighted loss \\
Occupancy \citep{mackenzie2002occupancy} & Presence and detection from repeated surveys & Same likelihood structure; no recommender use found \\
MRP, Fay-Herriot \citep{park2004mrp,fay1979smallarea} & Population and small-area estimates & Borrowed and tested on this target; weaker than DAX \\
NAVIP \citep{kim2022navip} & Propensity-weighted propagation in LightGCN & No lifelong data; no examination model \\
KuaiRand evaluations \citep{gao2022kuairand,zhang2026categorymeans} & Random insertion as the test set & No out-of-pool lifelong history in the test \\
\bottomrule
\end{tabular}
\end{table}

The novelty claim is therefore a combination and its controlled evidence, not any single ingredient. Each ingredient has
a precedent: lifelong behaviour statistics, examination multiplied by attractiveness, and Beta shrinkage towards a group
mean. As far as our searches went, none of the works in Table~\ref{tab:positioning} evaluates the combination in the form
used here, which has lifelong channels built from other pools, examination factors for the user and the scenario, and a
temporal protocol on randomly inserted items, with tuned baselines and split-level inference.
